\chapter{绪论}

\section{研究背景与意义}

随着载人航天任务复杂度的持续提升，航天员在交会对接、
姿态调整、轨道机动以及长期在轨操作过程中，需要面对更加复杂的
过载环境和运动耦合工况。为了在地面条件下尽可能真实地复现飞行
过程中的加速度刺激特性，构建具有较高动态响应能力、较好加速度
连续性以及较大工作空间的地面模拟装备，已经成为航天训练与飞行
器工程验证中的重要研究方向。早期的空间飞行加速度模拟设备多以
离心机为核心，其典型思想是通过大尺寸回转机构产生持续径向过载，
Barker 提出的空间飞行加速度模拟器即代表了这一技术路线的早期
探索\cite{barker1961simulator}。此后，围绕载人离心机的结构优化、
运动学与动力学分析、多目标构型设计以及抗荷效能评估等问题，
国内外学者开展了大量研究\cite{mohajer2021centrifugedesign,
winter2023centrifugeopt,tu2020gtolerance,liu2024centrifuge}。
这些研究表明，离心机系统在大幅值过载再现方面具有独特优势，但
其设备尺度大、建设与运行成本高、运动模式相对单一，且难以兼顾
多方向耦合加速度的灵活生成，因此在复杂机动环境的地面模拟方面
仍存在一定局限\cite{barker1961simulator,mohajer2021centrifugedesign,
winter2023centrifugeopt}。

除了传统离心机之外，基于并联机构或多自由度运动平台的地面运动
模拟方法也得到了广泛关注。相关研究表明，通过合理的机构设计与
控制策略，可以在有限空间内构造二维乃至三维动态加速度环境，从
而为特定任务场景下的训练与验证提供更具针对性的测试条件
\cite{zhang2020groundsimulation,liu2024centrifuge}。然而，这类系统
往往仍然面临动平台惯量偏大、结构复杂、驱动链条长以及高动态响
应能力受限等问题。当目标任务对“连续旋转”“大工作空间”“高加
速度响应”和“加速度平滑输出”提出同时要求时，传统刚性连杆机
构往往难以兼顾这些指标。因此，从结构形式上寻找更轻量、更具构
型灵活性的新型模拟平台，成为推动该类装备进一步发展的现实需求。

\begin{figure}[htbp]
  \centering
  \fbox{\begin{minipage}[c][5.0cm][c]{0.82\textwidth}
    \centering
    建议插入图：典型地面加速度模拟装备示意图。\\
    可将传统离心机、Stewart 平台、二维动态模拟平台等\\
    代表性方案放在同一张图中，体现技术演进脉络。
  \end{minipage}}
  \caption{典型地面加速度模拟装备示意}
  \label{fig:intro-training-equipment}
\end{figure}

绳驱并联机器人以柔性绳索替代部分或全部刚性连杆，兼具并联机构
受力路径短和索驱系统轻量化的特点，在大工作空间、高负载自重比、
低运动部件惯量以及快速动态响应等方面具有明显优势，因而成为近
年来并联机器人研究中的重要方向\cite{zhang2022stateoftheart,
wang2025review,chen2025cdprreview}。从工程应用上看，绳驱并联机
构已经被用于大型展演装备、物流搬运、建筑施工辅助、抓取操作以
及服务型机器人等场景，其中德国馆 EXPO 2015 项目的成功应用展
示了绳驱并联机器人在大尺度空间内进行高效、平滑、多任务运动控
制的工程可行性\cite{tempel2015expo}。与此同时，现代机器人在农业
喷雾、智能施药和精准作业等领域也呈现出对大范围运动、轻量化执
行与复杂环境适应能力的强烈需求\cite{lan2022sprayreview,
zhang2009greenhouse}。虽然这些应用场景与航天训练在任务目标上
不同，但它们共同指向了同一类技术需求，即在保证结构轻量化和控
制安全性的前提下，实现更大范围、更高速度和更强动态适应性的末
端运动输出。

对于面向加速度模拟任务的系统而言，绳驱并联机器人尤其适合承担
高动态运动生成问题。这是因为绳索只能承受拉力，系统天然具备冗
余驱动和张力重构的控制特性；在合理构型设计和张力分配策略支持
下，可通过各绳索张力协同作用，在动平台上形成所需的广义力和广
义力矩\cite{gouttefarde2007wrench,phan2021wrenchclosure}。对于需
要连续旋转输出的平面任务而言，平面四绳冗余驱动构型能够在二维
平移和绕平面法向旋转这三个自由度上实现较好的控制灵活性，其结
构简洁、理论分析相对清晰，也便于后续样机实现和算法验证
\cite{wang2022schonflies,gonzalez2023singleloop,
juarez2022dynamiccontrol}。因此，以平面四绳机构为基础，开展连
续旋转、高动态、面向加速度输出的绳驱并联机器人关键技术研究，
具有较强的理论意义和工程应用价值。

\begin{figure}[htbp]
  \centering
  \begin{tikzpicture}[x=1cm,y=1cm,>=latex]
    \tikzset{
      titlebox/.style={draw, rounded corners=2pt, fill=gray!12, minimum width=3.2cm, minimum height=0.85cm, align=center, font=\small\bfseries},
      metric/.style={draw, rounded corners=2pt, fill=gray!8, minimum width=2.65cm, minimum height=0.72cm, align=center, font=\small},
      scorebox/.style={draw, rounded corners=2pt, minimum width=2.25cm, minimum height=0.72cm, align=center, font=\small},
      iconline/.style={line width=0.9pt},
      good/.style={fill=green!55!black},
      medium/.style={fill=orange!80!black},
      weak/.style={fill=gray!55}
    }

    \node[titlebox] (c1) at (0,0) {传统离心机};
    \node[titlebox] (c2) at (4.4,0) {多自由度平台};
    \node[titlebox] (c3) at (8.8,0) {绳驱并联机器人};

    \begin{scope}[shift={(-0.2,-1.25)}]
      \draw[iconline] (0,0) circle (0.38);
      \draw[iconline] (0,0) -- (0.78,0.38);
      \draw[iconline] (0.78,0.38) circle (0.10);
      \draw[iconline,->] (0.18,0.52) arc[start angle=70,end angle=-220,radius=0.55];
    \end{scope}
    \begin{scope}[shift={(4.2,-1.25)}]
      \draw[iconline] (-0.75,-0.28) -- (0.75,-0.28);
      \draw[iconline] (-0.55,-0.28) -- (-0.22,0.26);
      \draw[iconline] (0.55,-0.28) -- (0.22,0.26);
      \draw[iconline] (-0.22,0.26) -- (0.22,0.26);
      \draw[iconline] (-0.22,0.26) -- (-0.48,0.66);
      \draw[iconline] (0.22,0.26) -- (0.48,0.66);
      \draw[iconline] (-0.48,0.66) -- (0.48,0.66);
    \end{scope}
    \begin{scope}[shift={(8.7,-1.25)}]
      \draw[iconline] (-0.72,0.55) rectangle (0.72,-0.55);
      \draw[iconline, fill=white] (0,0) rectangle (0.36,0.26);
      \foreach \x/\y in {-0.72/0.55, 0.72/0.55, -0.72/-0.55, 0.72/-0.55} {
        \draw[iconline] (\x,\y) -- (0.18,0.13);
      }
    \end{scope}

    \node[metric] (m1) at (-3.25,-2.3) {工作空间};
    \node[metric] (m2) at (-3.25,-3.15) {动平台惯量};
    \node[metric] (m3) at (-3.25,-4.0) {连续旋转能力};
    \node[metric] (m4) at (-3.25,-4.85) {动态响应};
    \node[metric] (m5) at (-3.25,-5.7) {加速度连续性};
    \node[metric] (m6) at (-3.25,-6.55) {建设成本};

    \newcommand{\scorecell}[4]{%
      \node[scorebox] at (#1,#2) {#3\hspace{0.35em}#4};}

    \scorecell{0}{-2.3}{大}{\tikz{\fill[good] (0,0) circle (2.2pt);}}
    \scorecell{4.4}{-2.3}{中}{\tikz{\fill[medium] (0,0) circle (2.2pt);}}
    \scorecell{8.8}{-2.3}{大}{\tikz{\fill[good] (0,0) circle (2.2pt);}}

    \scorecell{0}{-3.15}{低}{\tikz{\fill[good] (0,0) circle (2.2pt);}}
    \scorecell{4.4}{-3.15}{高}{\tikz{\fill[weak] (0,0) circle (2.2pt);}}
    \scorecell{8.8}{-3.15}{低}{\tikz{\fill[good] (0,0) circle (2.2pt);}}

    \scorecell{0}{-4.0}{强}{\tikz{\fill[good] (0,0) circle (2.2pt);}}
    \scorecell{4.4}{-4.0}{弱}{\tikz{\fill[weak] (0,0) circle (2.2pt);}}
    \scorecell{8.8}{-4.0}{强}{\tikz{\fill[good] (0,0) circle (2.2pt);}}

    \scorecell{0}{-4.85}{中}{\tikz{\fill[medium] (0,0) circle (2.2pt);}}
    \scorecell{4.4}{-4.85}{中}{\tikz{\fill[medium] (0,0) circle (2.2pt);}}
    \scorecell{8.8}{-4.85}{强}{\tikz{\fill[good] (0,0) circle (2.2pt);}}

    \scorecell{0}{-5.7}{中}{\tikz{\fill[medium] (0,0) circle (2.2pt);}}
    \scorecell{4.4}{-5.7}{中}{\tikz{\fill[medium] (0,0) circle (2.2pt);}}
    \scorecell{8.8}{-5.7}{强}{\tikz{\fill[good] (0,0) circle (2.2pt);}}

    \scorecell{0}{-6.55}{高}{\tikz{\fill[weak] (0,0) circle (2.2pt);}}
    \scorecell{4.4}{-6.55}{高}{\tikz{\fill[weak] (0,0) circle (2.2pt);}}
    \scorecell{8.8}{-6.55}{中}{\tikz{\fill[medium] (0,0) circle (2.2pt);}}

    \node[font=\scriptsize, anchor=west] at (-3.25,-7.45) {\tikz{\fill[good] (0,0) circle (2.2pt);} 优势 \qquad \tikz{\fill[medium] (0,0) circle (2.2pt);} 中等 \qquad \tikz{\fill[weak] (0,0) circle (2.2pt);} 相对较弱};
  \end{tikzpicture}
  \caption{典型加速度模拟技术路线对比示意}
  \label{fig:intro-compare-solutions}
\end{figure}

进一步地，从控制目标来看，若仅将绳驱并联机器人视作传统位置伺
服对象，则其轨迹规划与逆运动学求解通常以绳长或电机位置为主导
变量，能够较好地完成位置跟踪问题\cite{tang2015trajectory,
feiler2023mechatronic}。但对于加速度模拟任务而言，用户真正关注
的是末端平台所输出的加速度连续性、动态响应和受力品质，而不仅
是位置误差本身。已有研究表明，在绳索柔性、摩擦、离散轨迹求解
误差及机构耦合效应的共同影响下，单纯位置控制往往难以保证高动
态工况下的加速度平滑性，进而可能导致末端动态输出存在抖动甚至
局部不连续现象\cite{chen2025cdprreview,wang2024dynamictension}。
因此，相比以位置为唯一目标的控制方法，从广义力生成角度出发的
力控制和张力控制思路更适合本课题的应用目标
\cite{piao2019indirectforce,cao2023realtimetension,
song2024realtimemethod}。这一控制目标的转变，实际上也是本文从
“位置控制验证”进一步走向“张力分配与力控制算法研究”的核心原
因。

从工程实现角度看，连续旋转、高动态与绳驱约束之间并非天然兼容。
一方面，连续旋转能力要求机构在较大角位移范围内保持稳定绕线、
合理包角与较好的传力特性；另一方面，高动态输出要求系统具备较
快的驱动响应、较稳定的张力重构能力以及较强的轨迹与受力协调能
力\cite{gonzalez2023singleloop,juarez2022dynamiccontrol,
wang2024dynamictension}。这就意味着，单纯从几何层面完成机构布
置还远远不够，还需要在运动学建模、动力学分析、张力分配与控制
算法层面建立完整的研究链条。正是在这一背景下，本文选择平面四
绳冗余驱动机构作为研究对象，围绕其结构设计、位置控制验证、力
控制思路、张力分配算法以及实验实现开展研究。

综上所述，本文工作的研究意义主要体现在以下几个方面。其一，从
理论层面看，针对连续旋转平面四绳机构建立较为完整的建模与控制
分析框架，有助于完善绳驱并联机器人在高动态任务场景下的研究体
系\cite{zhang2022stateoftheart,wang2025review}。其二，从方法层面
看，通过将控制目标从位置跟踪进一步拓展到广义力调节与张力分配
优化，可为提升末端加速度连续性提供新的实现路径
\cite{su2016cablemass,cao2023realtimetension,wang2024dynamictension}。
其三，从工程层面看，相关研究可为航天员训练、飞行模拟和复杂运
动环境再现设备的研制提供参考，也可为其他需要大工作空间和高动
态响应的索驱机器人系统设计提供借鉴\cite{tempel2015expo,
lan2022sprayreview,zhang2009greenhouse}。

\section{国内外研究现状}

\subsection{加速度模拟与地面训练装备研究现状}

围绕航天员训练和飞行动力学再现问题，国外对地面加速度模拟装备
的研究起步较早。早期研究主要着眼于通过机械回转装置再现持续过
载环境，典型工作包括 Barker 提出的空间飞行加速度模拟设备
\cite{barker1961simulator}。这类系统为后续离心机平台的发展奠定
了基础，其核心思想是利用机构几何与角速度分布，在人体或载荷位
置处产生可预测的加速度场。随着机电技术和控制理论的发展，研究
重点逐渐从“能否产生过载”转向“如何提升加速度环境的真实性、舒
适性和可控性”。近年来，Mohajer 等围绕高效构型离心机系统开展了
运动学与动力学分析，指出合理配置回转结构参数能够改善系统的运
动性能和载荷响应\cite{mohajer2021centrifugedesign}；Winter 等进一
步从多目标优化角度分析了离心机系统性能提升问题，说明地面过载
装备在结构层面仍有较大优化空间\cite{winter2023centrifugeopt}；
Tu 等则从生理响应与抗荷效能角度分析了人体在离心机环境下的耐受
特征，说明训练装备的动力学特性与人体感知之间存在紧密耦合关系
\cite{tu2020gtolerance}。总体而言，传统离心机路线已经形成了较为
成熟的理论与工程基础，但其结构尺寸、运动模式与多自由度任务适
配性仍制约着进一步应用拓展。

国内在该方向同样开展了针对载人离心机和动态加速度环境构建的持
续研究。Zhang 等针对二维动态加速度环境的地面模拟试验进行了研
究，给出了特定场景下动态加速度环境生成的实验方法与验证结果，
说明通过机构设计与控制策略配合，可以在地面条件下模拟较为复杂
的动态加速度过程\cite{zhang2020groundsimulation}。刘晨晖等进一步
针对三轴载人离心机的加速度过载模拟算法开展开发与验证工作，体
现出国内在高维加速度环境建模、算法实现及实验校验方面的持续推
进\cite{liu2024centrifuge}。这些成果表明，国内学者已经认识到地面
模拟装备不仅要能够产生目标加速度，还必须兼顾动态一致性、系统
安全性与人机适配性。相比之下，如何在相对紧凑的结构中生成更丰
富、更连续的加速度环境，仍是值得继续突破的问题。

\begin{figure}[htbp]
  \centering
  \fbox{\begin{minipage}[c][4.8cm][c]{0.82\textwidth}
    \centering
    建议插入图：加速度模拟装备研究脉络图。\\
    可采用时间轴形式，标出空间飞行加速度模拟器、\\
    载人离心机、二维动态加速度模拟平台等代表性成果。
  \end{minipage}}
  \caption{加速度模拟装备研究脉络示意}
  \label{fig:intro-acceleration-history}
\end{figure}

\subsection{绳驱并联机器人构型与应用研究现状}

在新型运动平台研究方面，绳驱并联机器人因其结构轻量化和大尺度
运动潜力而受到广泛关注。Zhang 等对绳驱并联机器人的理论与应用
进行了系统综述，概括了其在构型设计、工作空间分析、控制方法及
工程应用等方面的发展现状\cite{zhang2022stateoftheart}；Wang 等则
从设计、建模与控制三方面对绳驱并联机器人的最新进展进行了综述，
指出绳驱系统未来的重要发展方向包括高精度建模、高动态控制以及
面向复杂任务的多模态应用\cite{wang2025review}。国内的综述工作也
表明，随着柔索动力学、张力分配和自适应控制等研究的深入，绳驱
机器人正在从传统的大型低速应用逐步走向高动态、高精度和复杂交
互任务\cite{chen2025cdprreview}。

从工程应用角度看，绳驱并联机器人已经不再局限于实验室环境。德
国馆 EXPO 2015 项目是典型的大型应用案例之一，Tempel 等对其设
计与编程过程进行了总结，说明绳驱并联机器人在大尺度空间中的协
调控制、任务组织和系统集成方面已经具备较高的工程成熟度
\cite{tempel2015expo}。在面向作业任务的机器人研究中，农业喷雾和
智能施药机器人也体现出对大工作空间、轻量化执行机构与复杂环境
适应能力的强烈需求\cite{lan2022sprayreview,zhang2009greenhouse}。
虽然这类应用与本文的加速度模拟场景并不相同，但它们说明绳驱或
轻量化驱动结构在复杂环境中的广泛适用性，也从侧面验证了绳驱并
联技术路线的工程潜力。

在构型研究方面，平面机构由于自由度清晰、建模难度相对适中，常
被作为绳驱并联机器人理论研究与原理验证的重要对象。Feiler 等从
机电一体化设计角度对一类平面绳驱并联机器人进行了系统分析，表
明构型参数、驱动布置和执行机构选择会直接影响系统的控制性能与
工程实现难度\cite{feiler2023mechatronic}。对于需要连续旋转输出的场
景，平面构型又具备路径关系直观、实验实现成本较低和便于算法验
证等优势，因此也成为本文采用平面四绳机构的重要原因。

\begin{figure}[htbp]
  \centering
  \fbox{\begin{minipage}[c][5.0cm][c]{0.82\textwidth}
    \centering
    建议插入图：绳驱并联机器人代表性应用与构型综述图。\\
    上半部分可放大型工程应用、农业/服务场景示例，\\
    下半部分可放平面、空间、悬挂式等典型机构草图。
  \end{minipage}}
  \caption{绳驱并联机器人应用与构型示意}
  \label{fig:intro-cdpr-applications}
\end{figure}

\subsection{工作空间、动力学建模与连续旋转研究现状}

工作空间与受力可行性是绳驱并联机器人研究中的基础问题。由于绳
索仅能承受拉力，末端执行器在某一位姿下能否产生期望广义力，不
仅取决于几何可达性，还受到正张力约束和张力上下限约束的共同影
响。Gouttefarde 等提出了绳驱并联机构的可受力工作空间概念，为分
析机构在给定位姿下的广义力生成能力提供了经典方法
\cite{gouttefarde2007wrench}；Phan Gia Luan 和 Nguyen Truong
Thinh 则从绳驱并联机构的扳手闭包条件出发，进一步讨论了机构实现
稳定受力输出的必要条件，为后续张力分配和构型优化研究提供了理
论基础\cite{phan2021wrenchclosure}。这些成果说明，绳驱并联机器人
的性能分析不能仅停留在几何运动学层面，还必须将张力约束与广义
力生成能力纳入统一框架中加以讨论。

在机构构型与动力学建模方面，学者们已经围绕不同类型的平面、空
间与悬挂式绳驱系统开展了大量研究。Wang 等针对满足 Sch{\"o}nflies
运动的悬挂式绳驱并联机器人建立了运动学、动力学与工作空间分析
模型，展示了绳驱机构在特定自由度任务中的建模思路
\cite{wang2022schonflies}。Gonz{\'a}lez-Rodr{\'i}guez 等研究了单绳环式
平面绳驱并联机器人的动力学模型，说明在平面构型下，通过特殊的
缠绕与传动方式可形成具有连续运动能力的系统结构
\cite{gonzalez2023singleloop}；Ju{\'a}rez-P{\'e}rez 等进一步在大可受力
工作空间条件下讨论了新型平面绳驱并联机器人的动态控制问题
\cite{juarez2022dynamiccontrol}。这些工作表明，围绕平面绳驱系统开
展连续旋转与高动态控制研究，既具有明确的理论基础，也具备较好
的工程可实施性。

从轨迹规划角度看，Tang 等针对平面二自由度冗余驱动绳悬挂机器人
开展了动态轨迹规划研究，说明轨迹规划结果与张力分布、驱动力需
求和系统动态响应之间存在显著耦合关系\cite{tang2015trajectory}。
这意味着，对于面向连续旋转和高动态输出的绳驱平台而言，单纯从
“可达”角度进行规划仍然不足，还必须综合考虑动力学负载和张力
演化过程。对本文而言，这一点尤为关键，因为本文的目标并非仅实
现几何意义上的连续旋转，而是要在连续旋转过程中兼顾末端加速度
输出品质。

\subsection{张力分配与力控制研究现状}

对于冗余驱动绳驱并联机器人而言，控制问题的核心之一在于张力分
配。由于系统自由度通常小于绳索数，在给定末端位姿或期望广义力
的条件下，绳张力解往往并不唯一，同时还必须满足绳索正张力约束、
张力上下限约束以及驱动器能力限制。因此，如何在满足受力平衡的
同时获得具有较好稳定性、实时性和可实施性的张力解，一直是国内
外研究的重要方向\cite{gouttefarde2007wrench,phan2021wrenchclosure,
cao2023realtimetension,song2024realtimemethod}。

在动力学建模与张力优化方面，苏宇等考虑了绳索质量和惯性力对系
统动力学的影响，研究了绳牵引并联机器人的动力学建模与张力优化
分布问题，说明在高动态工况下，忽略绳索分布质量和惯性效应可能
导致张力解与真实系统响应之间出现偏差\cite{su2016cablemass}。陈彦
霖等的综述工作则指出，绳驱机器人的动力学建模正在从刚性、无质
量绳索的简化假设逐步转向考虑绳索弹性、阻尼和复杂耦合效应的更
精细模型\cite{chen2025cdprreview}。这意味着，对于以高动态加速度输
出为目标的系统，仅依赖静态几何关系开展控制分析是不够的，必须
将动力学与张力演化过程结合起来加以研究。

在控制策略方面，陈科举等基于系统刚度开展了绳驱并联机器人的力
/位混合控制研究，体现了通过位置与力协同调节提升系统输出品质的
思路\cite{chen2023hybrid}；李永泉等针对冗余驱动并联机器人开展了
动力学参数辨识及控制研究，说明在冗余驱动系统中，精确模型辨识
与控制器设计之间存在密切关系\cite{li2019spherical}。在国外研究中，
Piao 等通过力传感器测量与神经网络估计相结合的方式实现了绳驱并
联机器人的间接力控制，展示了力控制在实际绳驱系统中的可实施性
\cite{piao2019indirectforce}；Cheah 等则从自适应与被动性角度提出了
适用于不同姿态参数化的绳驱并联机器人位姿跟踪控制方法，为复杂
模型条件下的稳定控制提供了新的理论支撑\cite{cheah2024adaptive}。
这些工作说明，绳驱系统的控制研究已逐步从单纯的几何跟踪扩展到
动力学一致性、鲁棒性和力学约束并重的阶段。

围绕张力分配算法本身，国内外学者也提出了多种求解方法。Cao 等
给出了面向绳驱并联机器人的实时张力分配设计方法，强调在保证计
算实时性的同时兼顾张力约束与系统稳定性\cite{cao2023realtimetension}；
Song 等提出了新的实时张力分配方法，从算法层面提高了求解效率与
在线应用能力\cite{song2024realtimemethod}；Wang 等进一步将动力学
建模与张力分配优化结合起来，讨论了高动态场景下张力分配对系统
性能的影响\cite{wang2024dynamictension}。这些研究为本文后续围绕
平面四绳连续旋转机构开展张力分配算法分析、仿真和优化提供了重
要理论与方法参考。

\begin{figure}[htbp]
  \centering
  \fbox{\begin{minipage}[c][5.0cm][c]{0.82\textwidth}
    \centering
    建议插入图：绳驱并联机器人理论研究主线图。\\
    可以“构型设计—工作空间—动力学建模—张力分配—\\
    力控制—实验验证”为主线，做成流程示意图。
  \end{minipage}}
  \caption{绳驱并联机器人理论研究主线示意}
  \label{fig:intro-cdpr-theory}
\end{figure}

综合来看，国内外关于绳驱并联机器人及其张力分配、动力学建模和
力控制的研究已取得大量成果，但针对“连续旋转平面四绳机构”
“高动态加速度输出”“位置控制到力控制的控制目标转换”以及
“面向实际样机验证的张力分配算法优化”这一组合问题，现有研究
仍缺少一条从机构设计、数学建模、算法推导到实验实现的完整闭环
路线。因此，围绕高动态连续旋转平面绳驱并联机器人开展系统研究，
既能够承接已有理论成果，又具有明确的工程需求牵引和应用落点。

\section{现有研究存在的问题}
\section{本文研究内容与技术路线}
\section{论文结构安排}
\section{本章小结}


